%% ED26-METHOD-minimum-description-length, original Automatica1978 allocation.
\chapter{Rissanen's Structural and Numerical Parameter Arguments}
\label{sec:p15-rissanen-structural-and-numerical-inference}

Rissanen's 1978 proposal estimates a time-series model's integer structure
and real parameters by accounting for the description of both the data and
the fitted model. The selected presentation is the complete original
article, not the later universal-integer-prior or normalized-maximum-
likelihood theories \cite{Rissanen1978Modeling}. Its concrete calculation is
for scalar linear models with Gaussian innovations. The resulting score
uses approximations to coding cost and to the effect of finite precision;
those approximations must remain visible when the score is used.

The introduction contrasts this route with Akaike's MAICE expression
\(N\log r+2k\), where the fit term is minimized negative log likelihood
up to a constant and \(k\) counts parameters. Rissanen accepts the role
of Kullback discrepancy but questions deriving the structural penalty
from a bias adjustment. His alternative charges for actually describing
the structure and its fitted numbers. That historical objection is not a
proof that every bias-corrected selection rule fails.
He also distinguishes his earlier 1976 entropy-based proposals, reporting
their failure for autoregressive order estimation and seeking a more direct
description interpretation. Those earlier works and their guarantees are
not proved or silently adopted here. Increasing a model's flexibility can
improve its raw fit while adding parameters whose transmission has been
ignored; the source's structural-estimation objection concerns that
unpenalized use, not a theorem about every possible likelihood-based rule.

\section{What the entropy and coding arguments establish}
\label{sec:p15-rmdl-coding}

\begin{hbdefinition}{def:p15-rmdl-lengths}{Data, models, code units and ideal lengths}
The observed sequence is \(x=(x_1,\ldots,x_N)\), with \(N\geq1\).
An agreed decoder and representation language distinguish a structure
description from a numerical parameter description and a conditional data
description. A real numerical value requires a finite precision, not an
infinitely long expansion. Binary code lengths are measured in bits;
\(-\log_2 q\) is an ideal information length. Gaussian calculations
below use nats, with \(1\) bit equal to \(\ln2\) nats. An ideal real
length need not be an integer or the length of an actual prefix code.
\end{hbdefinition}

\begin{hbtheorem}{thm:p15-rmdl-gibbs}{proved=annex:p15.rmdl.gibbs}{The source Gibbs inequality with countable safeguards}
Let \((p_i),(q_i)\) be nonnegative masses on a finite or countable index
set with the same finite total \(M\). Zero-weight entropy and
cross-entropy terms are defined to be zero. Put
\(D(p,q)=\sum_{i:p_i>0}p_i\ln(p_i/q_i)\), taking a term with
\(q_i=0<p_i\) as \(+\infty\). Its negative part is summable,
\(D(p,q)\geq0\), and equality holds exactly when \(p=q\).
For probability masses with finite entropy
\(H_2(p)=-\sum_i p_i\log_2p_i\),
\[
 -\sum_i p_i\log_2q_i=H_2(p)+D(p,q)/\ln2.
\]
The right side may be infinite. If both entropies are infinite, equality
of infinite cross-entropies is not a uniqueness statement.
\end{hbtheorem}

The paper's Section~2 uses this inequality first for an independent-symbol
source. The true distribution minimizes expected ideal length, not the
length of every individual realized string. If prefix lengths \(l_i\)
satisfy \(W=\sum_i2^{-l_i}\leq1\), normalizing their weights gives
\(q_i=2^{-l_i}/W\); the expected length is
\(H_2(p)+D(p,q)/\ln2-\log_2W\). Hence unused coding mass cannot
improve the entropy lower bound. The paper refers to arithmetic coding to
approach fractional ideal lengths; it does not supply that coding algorithm
or prove a later minimax-regret theorem.

\begin{hbtheorem}{thm:p15-rmdl-types}{proved=annex:p15.rmdl.types}{The exact type-class size and its limiting role}
For an alphabet of fixed size \(a\geq1\) and counts
\(k_1+\cdots+k_a=N\), the number of strings with those counts is
\(S=N!/\prod_i k_i!\). For distinct arbitrary binary descriptions,
including the empty word, their least possible uniform mean length is
\[
 b-\frac{2^{b+1}-b-2}{S},\qquad b=\lfloor\log_2(S+1)\rfloor.
\]
It is at least \(\log_2S-3\). Prefix descriptions have mean at least
\(\log_2S\). A finite enumerative prefix description specifies the
counts and then a rank within their type, with length at most
\(\lceil a\log_2(N+1)\rceil+\lceil\log_2S\rceil\).
Furthermore
\[
 \frac{\log_2S}{N}
 =H_2(k_1/N,\ldots,k_a/N)+O_a(\log(N+1)/N).
\]
Thus frequency convergence gives the type-class limit \(H_2(p)\).
It does not make a dependent process's entropy rate equal to its one-symbol
entropy: a stationary ergodic random-phase alternating binary process has
one-symbol entropy \(1\) but block entropy \(1\) for every \(N\).
\end{hbtheorem}

This supplies the full finite counting argument behind the source's sketch.
Its finite assertion that one arbitrary code must have at least
\(\log_2S\) bits needs the prefix/fixed-block convention: for \(S=3\),
the empty word, \(0\), and \(1\) are an injective nonprefix assignment.
The asymptotic per-symbol conclusion survives the constant-size correction.
Under an independent source, strings of one type are equiprobable; frequency
convergence alone does not establish that premise for a dependent source.

\section{Structure, innovations and the cost of numerical precision}
\label{sec:p15-rmdl-structure}

The source first proposes an invertible dynamic representation
\(x_t=f_\alpha(x_{t-1},\ldots,x_{t-n},e_t,\ldots,e_{t-m})\),
with zero initial conditions. Its observed-input version also contains
past \(u\)'s; the article then suppresses that input to focus on the
coding question. Invertibility means that innovations can be recovered
recursively from the observations and earlier innovations. It does not
mean that the innovations are independent. A statistical model for them
is an additional part of the description.

Write \(\theta=(s,\xi)\), with integer structure \(s\), such as the
orders \((n,m)\), and real parameters \(\xi\), including distribution
parameters. If innovations have \(q\) coordinates, resolution
\(\varepsilon\) and modeled density \(p_\beta\), the source's
conditional-length approximation is
\[
 L(e\mid\theta)=-\sum_{t=1}^N\ln\bigl(\varepsilon^q p_\beta(e_t)\bigr),
 \qquad L(x,\theta)=L(e\mid\theta)+L(\theta).
\]
This is a small-cell density approximation, not the exact cell probability.
The positive sign printed in source equation~(3.4) is repaired here:
negative log probability agrees with its Gibbs discussion and Gaussian
equation~(4.3). A density alone is not a dimensionless event probability.

The article approximates the parameter part by
\[
 L(\theta)\approx\ln\bigl((n+1)(m+1)\bigr)
                  +\sum_i\ln\frac{\lvert\xi_i\rvert}{\delta_i},
\]
with \(\delta_i\) the parameter's resolution. The absolute-value bars
here denote numerical magnitude, not cardinality. The source explicitly
omits sign costs and acknowledges a zero-indicator cost when a small
coordinate is omitted. Its reconstruction/code algorithms are omitted as
fixed overhead; they can be cancelled only when genuinely common to the
candidate descriptions. A larger model whose added parameter costs a
positive amount without reducing the data part has a larger total length.
That observation motivates nonredundant parameterizations, not a theorem
that every model representation has one canonical global optimum.

\begin{hbwarning}{def:p15-rmdl-code-boundary}{The approximate costs are not an unrestricted universal code}
The source's analogy assigns an integer \(j\) weight \(1/j\).
These weights do not normalize on all positive integers: each block
\(2^h\leq j<2^{h+1}\) has sum at least \(1/2\). On a declared
finite range \(1,\ldots,J\), normalized weights are
\(1/(j\sum_{h=1}^J1/h)\); the extra normalizer matters if the range
changes. Likewise the structure weights
\(1/((n+1)(m+1))\) do not form a probability on unbounded orders.
Signs, delimiters, zero flags and decoder specifications must be supplied
for a literal code. This 1978 approximation is not the later universal
integer code or normalized maximum-likelihood construction.
\end{hbwarning}

The source also suggests coding the parameter for a new batch conditional
on previous fitted parameters, motivated by near-Gaussian estimator laws.
It explicitly does not develop that alternative. It is not a justified
Gaussian prior for the single-batch calculation that follows.

\begin{hbtheorem}{thm:p15-rmdl-dependence}{proved=annex:p15.rmdl.dependence}{The dependence identity and the needed comparison condition}
For finite-valued innovations with joint law \(P\),
\[
 D\bigl(P,\prod_tP_t\bigr)/\ln2
       =\sum_tH_2(P_t)-H_2(P)\geq0,
\]
with equality exactly for independence. If each candidate innovation
transformation is a bijection on the same finite observation space,
\(H_2(P)\) is the same for every candidate; minimizing the marginal-
entropy sum then minimizes dependence. More generally, if their joint
entropies differ from one common value by at most \(\eta\), a minimizer
of the marginal-entropy sum has dependence within \(2\eta\) of any
candidate's dependence. This is an additive comparison, not mere convergence
of entropy ratios to one.
\end{hbtheorem}

This is the exact finite counterpart of source equations~(3.5)--(3.6).
Invertibility before rounding does not ensure a bijection afterwards.
For instance, scaling \(\{0,1,2,3\}\) by a sufficiently small positive
factor is invertible, but rounding to unit bins can collapse all four
points. Their observed distinctions cannot then be recovered at the
original resolution. Nor does a relative entropy error tending to zero
force its absolute error to vanish when the common entropy grows.
The source's small-resolution and almost-sure limiting discussion therefore
requires suitable quantization and stochastic convergence assumptions;
Gibbs' inequality by itself does not establish its general minimizer limit.

\section{The original scalar Gaussian calculation}
\label{sec:p15-rmdl-gaussian}

\begin{hbdefinition}{def:p15-rmdl-arma}{The source recurrence and length surrogate}
For nonnegative orders \(n,m\), set \(p=n+m\) and use
\[
 x_t+\sum_{j=1}^na_jx_{t-j}
      =e_t+\sum_{j=1}^mb_je_{t-j},\qquad x_t=e_t=0\ (t\leq0).
\]
The numerical vector is \(\xi=(r,\eta)\), where
\(\eta=(a_1,\ldots,a_n,b_1,\ldots,b_m)\) and \(r>0\) is innovation
variance. The recurrence defines \(e_t(\eta)\) and
\(E_N(\eta)=\sum_te_t(\eta)^2\). Put
\(\kappa_{n,m}=\ln((n+1)(m+1))\) and
\(K_N=N\ln(2\pi/\varepsilon^2)/2\). On an active-coordinate region
with every \(\eta_i\ne0\), the source's Gaussian length approximation is
\[
 L_\delta(r,\eta)=K_N+\frac N2\ln r+
 \frac{E_N(\eta)}{2r}+\ln\frac r{\delta_0}
       +\sum_{i=1}^p\ln\frac{\lvert\eta_i\rvert}{\delta_i}+\kappa_{n,m}.
\]
All precisions are positive. The formula is a real-valued local surrogate,
not a claim that these are integer lengths of a completed prefix code.
\end{hbdefinition}

\begin{hbtheorem}{thm:p15-rmdl-variance}{proved=annex:p15.rmdl.variance}{The numerical variance estimate is not ordinary ML}
For fixed \(\eta,\delta\) with \(E_N(\eta)>0\), the unique minimizer
over \(r>0\) is
\(\rho(\eta)=E_N(\eta)/(N+2)\). Its profiled value is
\[
 K_N+\frac{N+2}{2}(1+\ln\rho(\eta))-\ln\delta_0
        +\sum_{i=1}^p\ln\frac{\lvert\eta_i\rvert}{\delta_i}+\kappa_{n,m}.
\]
If \(E_N=0\), there is no positive minimizer: the surrogate tends to
\(-\infty\) as \(r\downarrow0\). A zero-variance Gaussian or a
boundary-constrained precision regime needs a different treatment.
\end{hbtheorem}

The \(+2\) comes from transmitting \(r\): ordinary Gaussian ML without
that term gives \(E_N/N\). Source~(4.6) differs from this direct substitution
into~(4.4) by a printed, model-independent \(-1/2\); that shift does not
change its minimizer but is not an exact equality with~(4.4).

\begin{hbtheorem}{thm:p15-rmdl-precision}{proved=annex:p15.rmdl.precision}{What the rounding calculation actually minimizes}
Let \(\Phi:O\to\mathbb R\) be \(C^3\) on an open
\(O\subseteq\mathbb R^d\), at a stationary fit \(\widehat\xi\).
Assume its rounding box
\(\widehat\xi+\prod_i[-\delta_i/2,\delta_i/2]\) lies in \(O\),
the third derivative norm there is bounded by \(M\), and the Hessian
diagonal entries \(h_i\) at the fit are positive. Independent errors
\(A_i\) uniform on \([-\delta_i/2,\delta_i/2]\) give
\[
 \mathbb E\Phi(\widehat\xi+A)=\Phi(\widehat\xi)
       +\sum_i h_i\delta_i^2/24+R,
 \qquad |R|\leq M(\sum_i\delta_i^2)^{3/2}/48.
\]
For the leading surrogate adding precision cost \(-\sum_i\ln\delta_i\),
the unique unconstrained minimizing widths satisfy
\(\delta_i^2=12/h_i\). This minimizes the declared quadratic surrogate,
not the exact expectation including \(R\).

In the paper's curvature approximation put \(d=p+1\),
\(G_i=\partial_{\eta_i}^2\ln\rho\) at the fitted \(\widehat\eta\),
\(h_0=(N+2)/(2\rho^2)\), and \(h_i=(N+2)G_i/2\).
If every \(G_i>0\), the source widths are
\[
 \delta_0=\rho\sqrt{24/(N+2)},\qquad
 \delta_i=\sqrt{24/((N+2)G_i)}.
\]
Define the published score
\[
 U_R=N\ln\rho+\sum_{i=1}^p\ln(\widehat\eta_i^2G_i)
                +d\ln(N+2)+2\kappa_{n,m}.
\]
Here the source-approximated rounding objective is explicitly
\[
 Q(\delta)=L_\delta(\rho,\widehat\eta)+\frac{N+2}{48}
       \left(\frac{\delta_0^2}{\rho^2}+\sum_{i=1}^pG_i\delta_i^2\right).
\]
The optimized quadratic surrogate \(Q_*\) instead satisfies exactly
\[
 2Q_*=U_R+2K_N+(N+2)+d(1-\ln24).
\]
Thus suppressing the last \(d\)-dependent term is not subtraction of
one common constant when different dimensions are compared.
\end{hbtheorem}

The article models truncation error as uniform and uses leading likelihood
curvature in~(4.7)--(4.9). Deterministic rounding does not automatically
have that error law, and the derivatives of magnitude-code costs need not
be negligible. The theorem states the approximation being optimized rather
than calling all the source's Hessian replacements exact identities.
The widths must also remain within the region supporting the Taylor bound.
If a \(G_i\) is zero or negative, there is no finite positive width from
that formula. For the variance width to satisfy \(\rho/\delta_0>1\)
requires \(N>22\); small samples need a boundary/actual-code analysis.

Source~(4.10) sets a coordinate to zero when
\(\lvert\widehat\eta_i\rvert/\delta_i<1\), removes its magnitude term
and reduces the active parameter count. The omitted zero flag still costs
something. A changed coordinate changes the recurrence and residual sum;
recomputing those is necessary before treating the reduced description as
a fitted alternative. This is not a universal significance test.

\begin{hbmethod}{def:p15-rmdl-procedure}{A defensible use of the source criterion}
Fix the observation resolution, decoder, candidate orders and admissible
numerical regions. For each order, compute innovations by the recurrence,
profile the variance, and fit the remaining parameters using the declared
length surrogate. Check positive residual sum, an interior active fit,
positive curvature and a valid small rounding region. Compute the widths,
check zero/precision branches, and recompute changed descriptions. Compare
the published score if that historical approximation is intended; compare
the constant-retaining surrogate if that is the intended objective.
For a literal finite code search, enumerate the specified finite candidate
messages and choose a shortest one. That finite minimum exists; it is not
an algorithmic guarantee of a global continuous or unrestricted structural
minimum. The paper also suggests joint minimization of its score, which
need not coincide with its staged, locally derived fitting calculation.
\end{hbmethod}

\section{Representation, BIC and the different Occam guarantee}
\label{sec:p15-rmdl-comparisons}

Coordinatewise nonzero scaling \(\psi_i=c_i\eta_i\) preserves
\(\eta_i^2G_i\), since \(G_i\) scales by \(c_i^{-2}\), and scales
the width by \(\lvert c_i\rvert\). Translation does not preserve the
magnitude code: \(\psi_i=\eta_i+b\) keeps the second derivative under
a linear translation but changes the product to \((\eta_i+b)^2G_i\).
The zero branch changes too. This is a concrete dependence on representation,
not a coordinate-free Fisher-volume or universal-coding theorem.

\begin{hblemma}{thm:p15-rmdl-leading}{proved=annex:p15.rmdl.leading}{A conditional leading penalty, not an MDL–BIC identity}
For a fixed order and active dimension, suppose along a sequence of fits
\(E_N/N\to v>0\), \(\widehat\eta_i\to c_i\ne0\), and
\(G_i\to g_i>0\). Then
\[
 U_R=N\ln(E_N/N)+d\ln N
      +\sum_i\ln(c_i^2g_i)+2\kappa_{n,m}-2+o(1).
\]
No consistency theorem for structural selection, varying dimension or
unrestricted model classes is asserted by this algebraic limit.
\end{hblemma}

The first and dimension terms resemble Schwarz's leading score, as the
original article notes. They do not make the objectives identical.
The Bayesian result at Theorem~\ref{thm:p15-bic-leading-order} starts from
proper prior densities and uniform Laplace conditions, retains different
constant-order terms, and concerns integrated evidence. This source instead
prices a transmitted fit with a sensitivity-dependent resolution. Neither
its \(1/j\) analogy nor a matching \(d\ln N\) term proves a posterior
truth probability. The entropy theorem likewise concerns a fixed outcome
description, not a justification for changing the coding language after
seeing the data.

\begin{hbtheorem}{thm:p15-rmdl-occam-case}{proved=annex:p15.rmdl.occam-case}{Exact short coding without the independent-sample guarantee}
Let \(X=\{0,1\}\), with uniform future input, true target \(f(x)=x\),
and fixed hypotheses \(h_0=0,h_1=1,f\). With \(N\geq1\) and the input
sequence known to the decoder, encode these hypotheses by \(0,10,110\),
respectively, and arbitrary label sequences by \(111\) followed by
\(N\) literal bits. This is a prefix code for a two-part description.
Selecting the shortest consistent description returns \(h_0\) on all-zero
labels, \(h_1\) on all-one labels, and \(f\) otherwise.
If all training inputs repeat one random uniform input, its future error
is exactly \(1/2\) with probability one, for every \(N\).
For genuinely independent uniform training inputs, the probability of that
error is instead \(2^{1-N}\). Realizability and a fixed representation
hold in both cases; independence is the premise that changes.
\end{hbtheorem}

This is an editorial comparison case, not a dataset or example claimed
printed by Rissanen. It separates a true minimum-code statement from the
uniform sampling guarantee of Theorem~\ref{thm:p15-occam-compression}.
The article's dynamic data may be dependent and noisy, whereas that PAC
theorem has specified independent inputs and noiseless target labels.
A short model/data description alone does not supply those hypotheses.
Rissanen's original conclusion is a structure-sensitive fitting proposal,
not a claim that compression identifies the true model on every sample.

\index{minimum description length!Rissanen1978}
\index{parameter precision!description length}
\index{structural estimation}
